\documentclass[11pt]{article}                 % Uses the article document class with 11-point font

% ============================================================
% Basic paper setup
% ============================================================
\usepackage[margin=1in]{geometry}             % Sets all page margins to 1 inch
\usepackage{setspace}                         % Provides commands for line spacing
\usepackage{amsmath,amssymb,amsthm}          % Provides advanced math, symbols, and theorem environments
\usepackage{mathtools}                        % Extends amsmath with additional mathematical tools
\usepackage{bm}                               % Provides bold mathematical symbols using \bm{}
\usepackage{booktabs}                         % Provides professional table rules such as \toprule
\usepackage{threeparttable}                   % Allows table notes below tabular environments
\usepackage{graphicx}                         % Allows insertion and resizing of graphics
\usepackage{subcaption}                       % Allows subfigures and subcaptions
\usepackage{float}                            % Provides additional float placement options such as [H]
\usepackage{placeins}                         % Provides \FloatBarrier to control float placement
\usepackage{enumitem}                         % Provides additional control over itemize and enumerate lists
\usepackage{natbib}                           % Provides citation commands such as \citet{} and \citep{}
\usepackage[hidelinks]{hyperref}              % Creates clickable references and links without colored boxes
\usepackage[nameinlink,noabbrev]{cleveref}    % Provides smart cross-references such as \cref{}

% ============================================================
% Formatting
% ============================================================
\onehalfspacing                               % Sets one-and-a-half line spacing throughout the paper
\setlength{\parindent}{1.5em}                 % Sets first-line paragraph indentation to 1.5 em
\setlength{\parskip}{0pt}                     % Removes extra vertical space between paragraphs

% ============================================================
% Common math shortcuts
% ============================================================
\newcommand{\E}{\mathbb{E}}                  % Expected value, e.g. \E[Y]
\newcommand{\Var}{\operatorname{Var}}        % Variance operator, e.g. \Var(Y)
\newcommand{\Cov}{\operatorname{Cov}}        % Covariance operator, e.g. \Cov(X,Y)
\newcommand{\R}{\mathbb{R}}                  % Set of real numbers, e.g. x \in \R
\newcommand{\ind}{\mathbf{1}}                % Indicator function, e.g. \ind\{X_i>0\}

% ============================================================
% Theorem-style environments
% ============================================================
\newtheorem{assumption}{Assumption}          % Creates numbered Assumption environments
\newtheorem{proposition}{Proposition}        % Creates numbered Proposition environments
\newtheorem{lemma}{Lemma}                    % Creates numbered Lemma environments

% ============================================================
% Title information
% ============================================================
\title{Paper Title}                           % Sets the title of the paper

\author{                                      % Begins author information
    Your Name\\                               % Author name; \\ creates a line break
    Department of Economics, University Name\\ % Department and university
    \texttt{your.email@university.edu}        % Displays the email address in typewriter font
}

\date{\today}                                 % Automatically inserts the compilation date

% ============================================================
% Document body
% ============================================================
\begin{document}                              % Begins the document body

\maketitle                                    % Prints the title, author, and date

% ============================================================
% Abstract
% ============================================================
\begin{abstract}                              % Begins the abstract environment
Write a concise abstract here. A typical research abstract briefly states the research question, data or setting, main method, main result, and contribution.

\vspace{0.5em}                                % Adds a small vertical space

\noindent\textbf{Keywords:} keyword 1; keyword 2; keyword 3
% \noindent removes paragraph indentation; \textbf{} makes text bold

\noindent\textbf{JEL Codes:} A10, C10
% Lists the relevant Journal of Economic Literature classification codes

\end{abstract}                                % Ends the abstract environment

% ============================================================
\section{Introduction}\label{sec:introduction}
% \section{} creates a numbered section
% \label{} assigns a reference label to the section
% ============================================================

Explain the research question, why it matters, what you do, what you find, and how the paper contributes to the literature.

A citation can look like \citet{example2020} or \citep{example2020}.
% \citet{} produces a textual citation, e.g. Example (2020)
% \citep{} produces a parenthetical citation, e.g. (Example, 2020)

% ============================================================
\section{Institutional Background or Model}\label{sec:background}
% ============================================================

Define the setting, model, institutions, or conceptual framework.

A generic specification can be written as

\begin{equation}\label{eq:baseline}           % Begins a numbered equation and assigns it a label
    y_i = \alpha + \beta x_i + \gamma' z_i + \varepsilon_i,
    \qquad i=1,\ldots,n.
\end{equation}                                % Ends the equation environment

% y_i            = outcome variable for observation i
% x_i            = main explanatory or treatment variable
% z_i            = vector of control variables
% \alpha         = intercept or constant term
% \beta          = coefficient of primary interest
% \gamma         = vector of coefficients on the controls
% \gamma' z_i    = inner product of \gamma and z_i
% \varepsilon_i  = error term
% i              = observation index
% n              = number of observations
% \qquad         = inserts horizontal mathematical spacing
% \ldots         = produces an ellipsis

Refer back to it as \cref{eq:baseline}.
% \cref{} automatically prints the object type, e.g. "equation (1)"

% ============================================================
\section{Data}\label{sec:data}
% ============================================================

Describe the data source, sample construction, key variables, and summary statistics.

\begin{table}[htbp}                           % Begins a floating table
% h = approximately here
% t = top of page
% b = bottom of page
% p = separate float page

\centering                                    % Horizontally centers the table

\caption{Summary Statistics}                  % Creates the table caption
\label{tab:summary}                           % Assigns a reference label to the table

\begin{threeparttable}                        % Allows table body plus notes

\begin{tabular}{lccc}                         % Creates a table with one left-aligned and three centered columns

\toprule                                      % Creates the top horizontal rule

Variable & Mean & Std. Dev. & Observations \\ % & separates columns; \\ ends the row

\midrule                                      % Creates the middle horizontal rule

Outcome & 1.25 & 0.42 & 1,000 \\              % First data row
Treatment & 0.48 & 0.50 & 1,000 \\            % Second data row
Control variable & 2.10 & 1.13 & 1,000 \\     % Third data row

\bottomrule                                   % Creates the bottom horizontal rule

\end{tabular}                                 % Ends the tabular environment

\begin{tablenotes}[flushleft]                 % Begins left-aligned table notes

\footnotesize                                 % Makes the table notes smaller

\item \textit{Notes:} This table reports summary statistics for the analysis sample.
% \item begins the note; \textit{} italicizes "Notes:"

\end{tablenotes}                              % Ends the table notes
\end{threeparttable}                          % Ends the three-part table
\end{table}                                   % Ends the table environment

% ============================================================
\section{Empirical Strategy}\label{sec:strategy}
% ============================================================

Explain the identifying variation and the assumptions required for interpretation. Keep the prose primary; equations should support the argument rather than replace it.

\begin{assumption}\label{assump:example}      % Begins a numbered assumption and assigns it a label
State an assumption clearly and in words that a reader can interpret.
\end{assumption}                              % Ends the assumption environment

% ============================================================
\section{Results}\label{sec:results}
% ============================================================

Present the main findings. Refer to tables and figures automatically, for example \cref{tab:summary}.
% \cref{tab:summary} automatically prints something like "Table 1"

\begin{figure}[htbp]                          % Begins a floating figure
\centering                                    % Centers the figure

% Replace the next box with:
% \includegraphics[width=0.75\linewidth]{figure_name.pdf}
% \includegraphics{} inserts an external figure
% width=0.75\linewidth scales it to 75 percent of the text width

\fbox{\rule{0pt}{2in}\rule{0.75\linewidth}{0pt}}
% \fbox{} draws a visible rectangular border
% \rule{0pt}{2in} sets the box height
% \rule{0.75\linewidth}{0pt} sets the box width

\caption{Main Result}                         % Creates the figure caption
\label{fig:main}                              % Assigns a reference label to the figure

\begin{minipage}{0.85\linewidth}              % Creates a text block 85 percent of the line width

\footnotesize                                 % Makes the figure notes smaller

\textit{Notes:} Explain what is plotted, the sample, uncertainty intervals, and any important construction details.
% \textit{} italicizes the word "Notes:"

\end{minipage}                                % Ends the minipage

\end{figure}                                  % Ends the figure environment

% ============================================================
\section{Robustness and Extensions}\label{sec:robustness}
% ============================================================

Organize robustness checks by the threat or concern they address, rather than listing specifications mechanically.

% ============================================================
\section{Conclusion}\label{sec:conclusion}
% ============================================================

Summarize the main result and contribution without simply repeating the introduction.

% ============================================================
% References
% ============================================================
\bibliographystyle{apalike}                   % Uses the apalike bibliography formatting style

% \bibliography{references}                   % Loads citations from references.bib
% Uncomment the command above when using a separate BibTeX file

% ============================================================
% Temporary example bibliography
% ============================================================
\begin{thebibliography}{1}                    % Begins a manual bibliography

\bibitem[Example(2020)]{example2020}          % Creates a bibliography entry with citation key example2020
Example, A. (2020).
\newblock Example paper title.                % Starts a new bibliography line
\newblock \emph{Journal Name}, 1(1), 1--20.   % \emph{} italicizes the journal name; -- creates an en dash

\end{thebibliography}                         % Ends the manual bibliography

% ============================================================
% Appendix
% ============================================================
\clearpage                                    % Starts the appendix on a new page
\appendix                                     % Changes subsequent section numbering to appendix format

\section{Additional Results}\label{app:additional}
% Creates Appendix A and assigns it a reference label

Put additional derivations, tables, figures, proofs, variable definitions, or robustness results here.

\end{document}                                % Ends the document